% ===========================================================================
\section{Threat Model and Trust Tiers}
\label{sec:threat_model}
% ===========================================================================

\subsection{Setting and Actors}
\label{sec:actors}

We consider a multi-organization AM supply chain in which a \emph{print
bureau} executes jobs against designs certified earlier in the lifecycle,
and downstream consumers (inspectors, certifiers, buyers) rely on the
on-chain record.  The bridge path comprises:

\begin{itemize}
  \item \textbf{Printer.} The machine and its firmware, reporting hotend
        and bed temperatures and job state over serial.  Consumer printers
        such as the Ender-3 have no secure element and no firmware
        attestation.
  \item \textbf{Bridge host.} An edge device (Raspberry Pi) running the
        print server, the telemetry recorder, and---at T1 and above---the
        device signing key.
  \item \textbf{Print bureau organization.} The Fabric member (MSP)
        operating the printer; \emph{economically motivated to appear
        reliable}, since its on-chain reputation gates future
        work~\cite{anonymous2026unified}.
  \item \textbf{Network.} The bridge-host-to-peer path (TLS-protected
        gRPC in our deployment).
\end{itemize}

The Fabric layer (ordering, endorsement, validation) is assumed correct;
consensus attacks are out of scope.  The adversary of interest is an
insider or intruder on the bridge path who wants the ledger to show a
conforming print that did not occur as recorded.

\subsection{Tamper Taxonomy}
\label{sec:taxonomy}

Table~\ref{tab:taxonomy} classifies the manipulations the bridge must
consider, ordered by where in the pipeline they occur.

\begin{table}[t]
\caption{Tamper taxonomy for the physical-to-chain bridge.}
\label{tab:taxonomy}
\centering
\scriptsize
\begin{tabular}{@{}p{2.3cm}p{3.0cm}p{1.9cm}@{}}
\toprule
\textbf{Manipulation} & \textbf{Description} & \textbf{Detected by} \\
\midrule
Post-hoc evidence tampering & Stored telemetry edited after submission to
hide a deviation & T1: ledger-hash mismatch \emph{and} signature failure \\
\addlinespace
Replay across assets & Evidence from a good print resubmitted for a
different part or design & T1: G-code SHA-256 binding; chaincode lifecycle
state machine \\
\addlinespace
Design substitution & Job executed from a G-code file other than the
certified one & T1: \icode{compliance} dimension scores 0 on hash mismatch \\
\addlinespace
Fabricated telemetry & Recorder fed false samples, or bundle synthesized
outright, then \emph{validly signed} & \textbf{Not} detected at T1
(trust boundary); T3 required \\
\addlinespace
Stolen / cloned device key & Signing key exfiltrated from the bridge host
and used elsewhere & T2 registry limits scope; revocation; T3 binds key to
hardware \\
\addlinespace
Colluding organization & Bureau admin controls host, key, and MSP identity
& T3 plus cross-organization inspection events \\
\bottomrule
\end{tabular}
\end{table}

\subsection{Trust Tiers}
\label{sec:tiers}

We organize bridge designs into four cumulative tiers
(Table~\ref{tab:tiers}).  Each tier strictly extends the guarantees of the
previous one; equally important is what each tier does \emph{not}
guarantee.

\begin{table*}[t]
\caption{Trust-tier ladder for the physical-to-chain bridge.  Each tier is cumulative over the previous one.}
\label{tab:tiers}
\centering
\scriptsize
\begin{tabular}{@{}lp{3.9cm}p{4.9cm}p{4.9cm}@{}}
\toprule
\textbf{Tier} & \textbf{Mechanism} & \textbf{Defends against} & \textbf{Does NOT defend against} \\
\midrule
T0 & Na\"ive submission: client POSTs event fields under an org MSP
identity &
Nothing on the bridge path; only on-chain immutability after commit &
Any manipulation before or at submission; evidence unauthenticated and
unbound to the artifact \\
\addlinespace
T1 & Edge-signed canonical bundle: device-resident Ed25519 key signs the
payload; SHA-256 of signed bundle committed on-chain (\textbf{this
paper}) &
Post-hoc evidence tampering (dual detection); replay across assets/designs
(G-code hash binding); silent design substitution; in-flight modification &
Fabrication at source by a device-key holder; key theft from the
unhardened host; bureau--host collusion \\
\addlinespace
T2 & On-chain device registry: chaincode verifies the bundle signature at
endorsement time; key enrollment/revocation are ledger events &
Bundles signed by unregistered or revoked keys; bounds a stolen key's
validity window; auditable device identity &
Fabrication at source by the \emph{registered} key holder; the registry
authenticates the key, not the physics \\
\addlinespace
T3 & Hardware root of trust: key in TEE/secure
element~\cite{costan2016sgx}, measured boot, sensor-level signing or
threshold attestation across independent observers &
Compromised host OS; key extraction; recorder substitution; raises
fabrication cost to hardware-attack level &
Physical-world spoofing by a fully malicious organization (e.g.,
detached thermistors) \\
\bottomrule
\end{tabular}
\end{table*}

\subsubsection{The T1 trust boundary, stated plainly}
\label{sec:boundary}

T1's guarantee is \emph{tamper-evidence after signing}, not \emph{truth at
the source}.  An adversary who controls the bridge host---and therefore
the device key---can synthesize a plausible telemetry bundle, sign it, and
commit its hash; nothing in T1 or T2 distinguishes this from an honest
capture, because the signature attests to the key, not to the physics.  T1
is nevertheless the correct first rung for legacy fleets: it converts the
dominant practical threats (retroactive log editing, evidence reuse, quiet
design substitution---all cheap, deniable, and otherwise undetectable)
into cryptographically detectable events, with no hardware changes.  The
fabrication-at-source residual is the gap T3 targets
(Section~\ref{sec:discussion}).

Two further properties come from the chaincode rather than the bridge.
The lifecycle state machine~\cite{anonymous2026unified} rejects a
\icode{PRINT\_JOB} event for an asset not in the
\icode{MATERIAL\_CERTIFIED} stage, so replayed evidence cannot attach to
an arbitrary asset even if its hash is fresh.  And because the bundle's
rating updates the operator's Beta posterior atomically with the
provenance event (Section~\ref{sec:arch_integration}), selectively
committing only flattering evidence leaves a sparse, low-confidence
rating history that the Wilson confidence-interval
gate~\cite{anonymous2026unified} penalizes directly.
